\newcommand{\DoNotLoadEpstopdf}{}
\pdfinfoomitdate=1
\pdftrailerid{}
\pdfsuppressptexinfo=-1
\documentclass[11pt,letterpaper]{article}
\usepackage[T1]{fontenc}
\usepackage{lmodern}
\usepackage{amsmath,amssymb,amsthm,mathtools}
\usepackage[margin=1in]{geometry}
\usepackage{booktabs,array,microtype,longtable}
\usepackage[hidelinks]{hyperref}
\hypersetup{pdftitle={Report 170 Labeled outerplanar graph expansions},pdfauthor={},pdfsubject={Exact coefficients component refinements and safe inversion for A097998 and A098000}}
\setlength{\parskip}{4pt}
\setlength{\parindent}{1.2em}
\setlength{\emergencystretch}{3em}
\setcounter{tocdepth}{1}
\numberwithin{equation}{section}
\newtheorem{theorem}{Theorem}[section]
\newtheorem{lemma}[theorem]{Lemma}
\newtheorem{proposition}[theorem]{Proposition}
\newtheorem{corollary}[theorem]{Corollary}
\theoremstyle{remark}\newtheorem{remark}[theorem]{Remark}
\newcommand{\dd}{\,\mathrm d}
\newcommand{\CC}{\mathbb C}
\newcommand{\QQ}{\mathbb Q}
\newcommand{\NN}{\mathbb N}
\newcommand{\EE}{\mathbb E}
\newcommand{\Rea}{\operatorname{Re}}
\newcommand{\Var}{\operatorname{Var}}
\newcommand{\Pois}{\operatorname{Poisson}}
\newcommand{\TV}{d_{\mathrm{TV}}}
\newcommand{\ceil}[1]{\left\lceil #1\right\rceil}
\title{Labeled outerplanar graph expansions\\\large Exact coefficients component refinements and safe inversion}
\author{Report 170}
\date{3 October 2026}
\begin{document}
\maketitle
\begin{abstract}
Let $c_n$ and $g_n$ count connected and arbitrary outerplanar simple
graphs on the fixed label set $[n]$, respectively. These are the
positive-index terms of OEIS A097998 and A098000. Starting with the
classical block decomposition, we give an explicit finite algebraic
rule for every fixed-order correction to
\[
 c_n\sim A n!\rho^{-n}n^{-5/2},\qquad
 g_n\sim e^\nu A n!\rho^{-n}n^{-5/2},
\]
and evaluate the first four corrections. An exact component-marked
coefficient identity leads to a uniform complex-parameter saddle proof
with a Gaussian-integrable remainder and no logarithmic loss.

For the component count $K_n$ we obtain second-order expansions of its
mean, variance, and connectivity probability, and prove
\[
 \TV\bigl(\mathcal L(K_n),\mathcal L(1+\Pois(\nu+5\tau/(2n)))\bigr)
 =O(n^{-2}).
\]
A logarithmic expansion gives a Lambert-$W$ starting point and a
rounding-safe, two-ceiling enclosure for the generalized inverse.
The companion supplies exact recurrences, algebraic certificates,
literal small-graph enumeration, independent formal coefficient
checks, and bounded regression tests.

The leading equivalents and limiting shifted-Poisson law are established
results. General subcritical graph theory already implies fixed-order
expansions. This report is an explicit specialization and reproducible
refinement, without a worldwide-priority claim. Decimal evaluations are
not certified intervals; inverse error constants and starting thresholds
are existential rather than evaluated.
\end{abstract}
\clearpage
\tableofcontents
\bigskip

\section{Objects conventions and scope}

An outerplanar graph is an abstract simple graph that admits a planar
embedding with all its vertices on the unbounded face. We count each
edge set on $[n]=\{1,\ldots,n\}$ once, regardless of the number of its
embeddings. Neither plane maps nor unlabelled isomorphism classes are
being counted.

For $n\geq1$, put $c_n=\mathrm{A097998}(n)$, and for $n\geq0$ put
$g_n=\mathrm{A098000}(n)$ \cite{oeisconn,oeisall}. The initial terms are
\[
\begin{array}{c|rrrrrr}
n&0&1&2&3&4&5\\\hline
c_n\text{ in the EGF}&0&1&1&4&37&602\\
g_n&1&1&2&8&63&893
\end{array}
\]
The OEIS connected sequence conventionally prepends $1$ at index zero.
That convention is not used in the connected exponential generating
function (EGF):
\begin{equation}\label{eq:egfs}
 C(z)=\sum_{n\geq1}c_n\frac{z^n}{n!},\qquad
 G(z)=\sum_{n\geq0}g_n\frac{z^n}{n!}=e^{C(z)}.
\end{equation}
Including the prepended term in $C$ would incorrectly multiply $G$ by
$e$. A block means a 2-connected outerplanar graph on at least three
vertices, or the single edge $K_2$. There is no singleton block.

Bodirsky and Kang develop exact labeled enumeration and sampling
\cite{bk}; Bodirsky, Gim\'enez, Kang, and Noy (BGKN) establish the leading
count equivalents and the limiting component law \cite{bgkn}. Kang's
thesis also supplies the correct analytic amplitude relation
\cite{kang}. The analytic framework for labeled subcritical graph
classes includes this family \cite{subcritical}. We retain these
attributions throughout. The useful content here is the finite exact
coefficient generator, four computed corrections, explicit finite-size
component refinements, a direct complete remainder proof, and the
integer-safe inverse construction.

Every asymptotic assertion below takes its order $J$ or $m$ fixed before
$n$ tends to infinity. No convergence of the resulting infinite formal
series is asserted. Uniformity in a marker means uniformity on each
fixed compact subset of $\CC$, not for markers growing with $n$.

\section{Exact block and marked identities}

\subsection{Block generating function}
Let $B$ be the block EGF and write $b=B'$. The classical polygon-dissection
normalization gives
\begin{equation}\label{eq:block}
 b(u)=\frac{1+5u-\sqrt{1-6u+u^2}}8,\qquad B(u)=\int_0^u b(t)\dd t.
\end{equation}
The square root is the branch equal to $1$ at zero. In particular,
\begin{equation}\label{eq:blockseries}
 b(u)=u+\frac{u^2}{2}+\frac{3u^3}{2}
          +\frac{11u^4}{2}+\frac{45u^5}{2}+O(u^6).
\end{equation}
For a 2-connected outerplanar graph with $n\geq3$, its outer Hamilton
cycle is unique. There are $(n-1)!/2$ unoriented Hamilton cycles on the
fixed labels, and for each cycle the allowable chords are polygon
dissections. Thus the dissection count is multiplied by $(n-1)!/2$,
not by the number of plane embeddings. The edge $K_2$ supplies the
linear term $u$ in $b$ separately \cite{bk}.

For an independent recurrence formulation, let
\[
 S(u)=1+uS(u)+uS(u)^2=1+2u+6u^2+22u^3+90u^4+\cdots.
\]
The branch with constant term $1$ is
$S=(1-u-\sqrt{1-6u+u^2})/(2u)$, and
$b=u+u(S-1)/4$. This also checks \eqref{eq:blockseries}. For example,
there are nine 2-connected outerplanar graphs on four labels, so the
coefficient of $u^3$ in $b$ is $9/3!=3/2$.

\subsection{Rooting and integration}
Define the vertex-rooted connected EGF $T=zC'(z)$. The rooted block
set decomposition and subsequent integration give
\begin{equation}\label{eq:rooted}
 T=z\exp(b(T)),\qquad C(z)=\psi(T(z)),\qquad
 \psi(u)=u-ub(u)+B(u).
\end{equation}
Indeed $\psi'(u)=1-ub'(u)$, and
$\dd z/z=(1/u-b'(u))\dd u$ under $z=u e^{-b(u)}$.
Both sides of $C=\psi(T)$ vanish at zero. Lagrange inversion applied
to $T$ therefore gives the exact connected coefficient identity
\begin{equation}\label{eq:connexact}
 \frac{c_n}{n!}=\frac1{n^2}[u^{n-1}]e^{nb(u)},\qquad n\geq1.
\end{equation}

The component-marked polynomial is
\begin{equation}\label{eq:fn}
 f_n(v)=n![z^n]e^{vC(z)}.
\end{equation}
For $n\geq1$, $f_n(1)=g_n$, $f_n(0)=0$, and $f_n'(0)=c_n$.
For a uniformly chosen graph counted by $g_n$, the probability generating
function (PGF) of its number $K_n$ of components is $f_n(v)/g_n$.

\begin{proposition}[Exact marked cancellation]\label{prop:marked}
For every $n\geq1$ and every complex $v$,
\begin{equation}\label{eq:marked}
 \frac{f_n(v)}{n!}=\frac{v}{n^2}[u^{n-1}]
 e^{v\psi(u)}\bigl(1+vu\psi'(u)\bigr)e^{nb(u)}.
\end{equation}
\end{proposition}
\begin{proof}
Lagrange inversion initially gives
$(v/n)[u^{n-1}]e^{v\psi}\psi'e^{nb}$. Set
$A_v(u)=e^{v\psi(u)}$ and $E(u)=e^{nb(u)}$. Coefficient differentiation
implies
\[
 (n-1)[u^{n-1}]A_vE
  =[u^{n-1}]u(A_v'+nb'A_v)E.
\]
Rearranging, and using $A_v'=v\psi'A_v$, gives
\[
 [u^{n-1}]A_vE(1-ub')
   =\frac1n[u^{n-1}]A_vE(1+vu\psi').
\]
Since $\psi'=1-ub'$, this proves \eqref{eq:marked}.
\end{proof}

The cancellation removes an amplitude that vanishes at the saddle.
Keeping its exact extra factor $1/n$ is essential for the unrooted
exponent $n^{-5/2}$ and its amplitude. No asymptotic approximation has
been used in these identities.

\section{Saddle constants and exact coefficient generator}

Let $R=3-2\sqrt2$, the positive branch singularity of $b$.
The nonnegative coefficients of $b$ show that $ub'(u)$ is strictly
increasing from zero to infinity on $(0,R)$. There is a unique
$\tau\in(0,R)$ with
\begin{equation}\label{eq:saddle}
 \tau b'(\tau)=1.
\end{equation}
Equivalently, it is the root in $(0.17076,0.17077)$ of
\begin{equation}\label{eq:quartic}
 3\tau^4-28\tau^3+70\tau^2-58\tau+8=0
\end{equation}
that satisfies the unsquared equation \eqref{eq:saddle}. The unsquared
condition fixes the radical branch; merely squaring an equation is not
a root-selection rule.

Put
\begin{equation}\label{eq:constants}
 \rho=\tau e^{-b(\tau)},\quad \nu=\psi(\tau),\quad
 D=u\frac{\dd}{\dd u},\quad \kappa_j=(D^jb)(\tau),\quad
 A=\frac{\tau}{\sqrt{2\pi\kappa_2}}.
\end{equation}
Then $\kappa_1=1$ and $\kappa_2>0$. Useful approximate values are
\[
\begin{aligned}
 \tau&=0.170764986812629023408139791998\ldots,\\
 \rho&=0.136593733696039583904892236367\ldots,\\
 \rho^{-1}&=7.32098005480462334884554918493\ldots,\\
 \nu&=0.148886742687358857557744764544\ldots,\\
 A&=0.00697609450712314901934080444574\ldots,\\
 e^\nu A&=0.00809604747608070499620425650041\ldots.
\end{aligned}
\]
These are numerical evaluations, not certified decimal intervals.
The exact definitions are \eqref{eq:saddle}--\eqref{eq:constants}.

At the selected root,
\[
 \sqrt{1-6\tau+\tau^2}=\frac{\tau(3-\tau)}{8-5\tau}.
\]
Consequently all the $\kappa_j$ belong to $\QQ(\tau)$. The correction
coefficients below also lie in this field. The leading constant $A$
itself includes the universal Gaussian factor $1/\sqrt{2\pi}$;
$A\sqrt{2\pi}=\tau/\sqrt{\kappa_2}$ is algebraic. We do not call $A$
an algebraic number.

Define formal analytic germs at $t=0$ by
\begin{align}
 K(t)&=b(\tau e^t)=\sum_{j\geq0}\frac{\kappa_j}{j!}t^j,\label{eq:K}\\
 Q(t)&=\psi(\tau e^t)-\nu,\qquad
 Q'(t)=\tau e^t(1-K'(t)),\quad Q(0)=0,\label{eq:Q}\\
 M_v(t)&=e^{t+vQ(t)}(1+vQ'(t)).\label{eq:M}
\end{align}
Here $Q(t)=O(t^2)$ and $M_v(0)=1$.
Let $\EE$ denote the formal Gaussian moment operation on polynomials
in $x$, determined by
$\EE x^{2h}=(2h-1)!!$ and $\EE x^{2h+1}=0$, with $(-1)!!=1$.

\begin{theorem}[All fixed orders with a complex marker]\label{thm:main}
For every fixed $J\geq0$ and every compact $\mathcal K\subset\CC$,
\begin{equation}\label{eq:main}
 \frac{f_n(v)}{v e^{v\nu}A n!\rho^{-n}n^{-5/2}}
 =\sum_{j=0}^{J}\frac{d_j(v)}{n^j}
       +O_{\mathcal K,J}(n^{-J-1}),
\end{equation}
uniformly for $v\in\mathcal K$, with the removable value used at $v=0$.
Here $d_0=1$, $d_j\in\QQ(\tau)[v]$, $\deg_v d_j\leq j$, and the exact
finite coefficient rule is
\begin{equation}\label{eq:generator}
 d_j(v)=[\epsilon^{2j}]\EE\left[
 M_v\left(\frac{i\epsilon x}{\sqrt{\kappa_2}}\right)
 \exp\left\{\sum_{k\geq3}
 \frac{\kappa_k(i x)^k\epsilon^{k-2}}
      {k!\kappa_2^{k/2}}\right\}\right].
\end{equation}
Only $k\leq2j+2$ and amplitude powers at most $2j$ are needed.
In particular,
\begin{align}
 c_n&=A n!\rho^{-n}n^{-5/2}
 \left(\sum_{j=0}^{J}\frac{d_j(0)}{n^j}+O_J(n^{-J-1})\right),\label{eq:conna}\\
 g_n&=e^\nu A n!\rho^{-n}n^{-5/2}
 \left(\sum_{j=0}^{J}\frac{d_j(1)}{n^j}+O_J(n^{-J-1})\right).\label{eq:alla}
\end{align}
\end{theorem}

For readers preferring a formula without a formal expectation, put
$a_\ell(v)=[t^\ell]M_v(t)$. Sum over $\ell\geq0$ and nonnegative
integers $m_k$, $3\leq k\leq2j+2$, satisfying
\[
 \ell+\sum_{k=3}^{2j+2}(k-2)m_k=2j.
\]
If $q=\ell+\sum k m_k$, then
\begin{equation}\label{eq:finitesum}
 d_j(v)=\sum
 \frac{a_\ell(v)(-1)^{q/2}(q-1)!!}{\kappa_2^{q/2}}
 \prod_{k=3}^{2j+2}\frac{(\kappa_k/k!)^{m_k}}{m_k!}.
\end{equation}
The constraint makes $q$ even. This is a finite algebraic formula for
every coefficient and makes no reference to fitted count data.

\section{Complete uniform saddle proof}

\subsection{Exact contour and the unique dominant point}
Cauchy's formula on $u=\tau e^{i\theta}$ transforms
Proposition~\ref{prop:marked} into the exact identity
\begin{equation}\label{eq:contour}
 f_n(v)=\frac{n!v\tau e^{v\nu}\rho^{-n}}{n^2}
 \frac1{2\pi}\int_{-\pi}^{\pi}M_v(i\theta)e^{n\Phi(\theta)}\dd\theta,
\end{equation}
where
\[
 \Phi(\theta)=K(i\theta)-K(0)-i\theta.
\]
Every power of $n$ and $\tau$ in \eqref{eq:contour} is exact. The extra
$e^{i\theta}$ from extracting $[u^{n-1}]$ is the factor $e^t$ in $M_v$.

Write $b(u)=\sum_{k\geq1}b_ku^k$. Since $b_k\geq0$ and $b_1=1$,
\begin{equation}\label{eq:global}
 \Rea\Phi(\theta)
 =\sum_{k\geq1}b_k\tau^k(\cos(k\theta)-1)
 \leq\tau(\cos\theta-1).
\end{equation}
On $[-\pi,\pi]$, this is strictly negative except at zero. Thus the
whole integration circle has one dominant saddle. On closed arcs
away from zero there is a fixed strictly negative bound.
Also
\[
 \Phi(\theta)=-\frac{\kappa_2\theta^2}{2}+O(\theta^3),
\]
so there are fixed $\eta,c>0$ with
$\Rea\Phi(\theta)\leq-c\theta^2$ for $|\theta|\leq\eta$.
The amplitude is uniformly bounded for $v$ in each fixed compact set
and $\theta\in[-\pi,\pi]$, because $\tau<R$ and its dependence on $v$
is entire. These observations prove both global dominance and marker
uniformity without invoking a $\Delta$-domain for $T(z)$.

\subsection{An integrable remainder}
Let $a=\sqrt{\kappa_2}$ and $\epsilon=n^{-1/2}$. On the central arc use
\[
 \theta=\epsilon x/a,\qquad |x|\leq\epsilon^{-1/5}.
\]
The angular cutoff is $\epsilon^{4/5}/a$, a fixed multiple of $n^{-2/5}$.
Set
\[
 R_3(t)=K(t)-K(0)-t-\kappa_2t^2/2=t^3h(t),
\]
where $h$ is analytic in a fixed neighborhood of zero, and define
\begin{align*}
 S(s,x)&=s^{-2}R_3(isx/a)=s(ix/a)^3h(isx/a),\\
 U_v(s,x)&=M_v(isx/a)e^{S(s,x)}.
\end{align*}
The value at $s=0$ is understood by analytic continuation. After leading
normalization, the central integral is
\begin{equation}\label{eq:central}
 \frac1{\sqrt{2\pi}}
 \int_{|x|\leq\epsilon^{-1/5}}e^{-x^2/2}U_v(\epsilon,x)\dd x.
\end{equation}

Fix $J$ and put $N=2J+2$. For $0\leq s\leq\epsilon$ on this arc,
the arguments $isx/a$ lie in a fixed smaller disk of analyticity once
$\epsilon$ is small. Uniform derivative bounds give
\[
 |S(s,x)|\leq C\epsilon|x|^3\leq C\epsilon^{2/5}\leq C,
 \qquad |\partial_s^r S(s,x)|\leq C_r|x|^{r+2}\quad(1\leq r\leq N).
\]
For completeness, writing $w=ix/a$ shows the latter estimate directly:
\[
 \partial_s^r S=s w^{r+3}h^{(r)}(sw)
                       +r w^{r+2}h^{(r-1)}(sw).
\]
The factor $s|w|$ is uniformly bounded. Similarly,
$|\partial_s^rM_v(isx/a)|\leq C_{\mathcal K,r}|x|^r$.
The product rule and the finite Bell-polynomial formula for derivatives
of an exponential now imply
\begin{equation}\label{eq:uderiv}
 |\partial_s^r U_v(s,x)|\leq C_{\mathcal K,r}(1+|x|)^{3r},
 \qquad 0\leq r\leq N.
\end{equation}
Indeed, a product of derivatives of $S$ whose derivative orders sum to
$r$ has degree at most $r+2r=3r$ in this bound; derivatives of $M_v$
do not increase it. The factor $e^S$ is bounded by $e^C$.

Taylor's theorem in the real parameter $s$, with its integral remainder,
therefore gives
\begin{align}
 U_v(\epsilon,x)&=\sum_{k=0}^{2J+1}p_k(v,x)\epsilon^k+E_n(v,x),
 \label{eq:taylor}\\
 |E_n(v,x)|&\leq C_{\mathcal K,J}\epsilon^{2J+2}(1+|x|)^{6J+6}.
 \label{eq:remainder}
\end{align}
Keeping the polynomial inside the Gaussian integral yields
\[
 \int e^{-x^2/2}|E_n(v,x)|\dd x
 \leq C'_{\mathcal K,J}\epsilon^{2J+2}.
\]
Thus there is no logarithmic loss. Taking the supremum of the polynomial
on a growing interval before integrating would not give the required
sharp estimate.

Each $p_k$ has parity $k$ in $x$: an amplitude monomial has equal
$x$- and $\epsilon$-degrees, while a phase monomial has degrees $q$
and $q-2$. All odd terms integrate to zero on the symmetric interval.
Extending each even polynomial integral to the real line incurs an
error smaller than every fixed power of $\epsilon$. The omitted
intermediate arc contributes at most a constant times
\[
 \epsilon^{-1}e^{-c'\epsilon^{-2/5}},
\]
and the distant arc contributes
$O_{\mathcal K}(\epsilon^{-1}e^{-c''/\epsilon^2})$ after leading
normalization. Both are smaller than every fixed power. Equations
\eqref{eq:central}--\eqref{eq:remainder} prove \eqref{eq:main}, with
\eqref{eq:generator} obtained by the finite Taylor multiplication.

\subsection{Algebraic coefficients and marker degree}
Since $Q'=\tau e^t(1-K')$ and $Q(0)=0$, all Taylor coefficients of $Q$
lie in $\QQ(\tau)$. Gaussian moments supply rational factors; only even
total powers of $x$ survive, so the square roots of $\kappa_2$ cancel.
Hence $d_j(v)\in\QQ(\tau)[v]$.

For $h\geq1$, the coefficient of $v^h$ in
$e^{vQ}(1+vQ')$ starts at degree at least $2h-1$ in $t$, because
$Q=O(t^2)$ and $Q'=O(t)$. Multiplication by $e^t$ does not lower that
degree. To contribute to $\epsilon^{2j}$ requires $2h-1\leq2j$,
which for integer $h$ gives $h\leq j$. This proves the degree bound.
The proof of Theorem~\ref{thm:main} is complete.

\section{Four corrections and independent coefficient checks}

Set $d=d_1(0)$ and define
\begin{align}
 d&=-\frac1{2\kappa_2}+\frac{\kappa_3}{2\kappa_2^2}
      +\frac{\kappa_4}{8\kappa_2^2}
      -\frac{5\kappa_3^2}{24\kappa_2^3},\label{eq:d1}\\
 t_2&=\tau\left(\frac1{\kappa_2}-\frac{\kappa_3}{3\kappa_2^2}\right).
 \label{eq:t2}
\end{align}
Then the first two complete marker polynomials are
\begin{align}
 d_1(v)&=d+\frac{5\tau}{2}v,\label{eq:d1v}\\
 d_2(v)&=d_2(0)+\left(\frac72\tau d-\frac{35}8t_2\right)v
                     +\frac{35}8\tau^2v^2.\label{eq:d2v}
\end{align}
Here \eqref{eq:generator} at $j=2,v=0$ is a short exact definition of
$d_2(0)$; no decimal is needed to specify the theorem. For example,
\[
 [t]M_v=1-v\tau\kappa_2,\qquad
 [t^2]M_v=\frac12-\frac{v\tau}{2}(5\kappa_2+\kappa_3).
\]
Substitution in \eqref{eq:finitesum} gives \eqref{eq:d1v} directly.
The analogous finite algebra gives \eqref{eq:d2v}.

\begin{center}
\begin{tabular}{crr}
\toprule
$j$&Connected $d_j(0)$&All $d_j(1)$\\
\midrule
1&2.28327991621488764390&2.71019238324646020242\\
2&5.05931216558214791991&7.37088006509066565629\\
3&12.62518674698345424780&22.53510992886156396083\\
4&37.15520265458457580487&78.66695831469578724957\\
\bottomrule
\end{tabular}
\end{center}
The table uses approximate evaluations of exact algebraic quantities.
For all four rows, the companion encodes the exact values in the basis
$1,\tau,\tau^2,\tau^3$ of
\[
 \QQ[t]/(3t^4-28t^3+70t^2-58t+8),
\]
with the real root selected by \eqref{eq:saddle}. As one readable example,
\begin{equation}\label{eq:fieldd}
 d=\frac{146018067}{47538652}
 -\frac{892831981}{178269945}\tau
 +\frac{834955523}{356539890}\tau^2
 -\frac{63474793}{237693260}\tau^3.
\end{equation}
The finite generator is a more compact exact presentation of higher
orders than expanding every large rational coefficient in print.

\subsection{An independent formal singular expansion}
This subsection checks coefficients independently; the analytic proof
is already complete. Put $X=\sqrt{1-z/\rho}$ and
$t=\log(T/\tau)$. Then
\[
 1-X^2=\exp(t-K(t)+K(0)),\qquad
 W(t)=\frac{1-\exp(t-K(t)+K(0))}{t^2},\quad W(0)=\kappa_2/2.
\]
The physical branch is $X=-t\sqrt{W(t)}$. Formal Lagrange inversion gives
\begin{equation}\label{eq:formalT}
 T_k:=[X^k]T
 =\frac{\tau(-1)^k}{k}[t^{k-1}]e^tW(t)^{-k/2},\qquad k\geq1,
\end{equation}
with $T_0=\tau$. Termwise integration of $zC'=T$ gives
\begin{equation}\label{eq:formalC}
 C_0=\nu,\quad C_1=0,\qquad
 C_j=-\frac2j\sum_{\substack{r\geq0\\j-2-2r\geq0}}T_{j-2-2r}
 \quad(j\geq2).
\end{equation}
In particular,
\[
 T_1=-\tau\sqrt{2/\kappa_2},\quad T_2=t_2,\quad C_2=-\tau,
 \quad C_3=\frac23\tau\sqrt{2/\kappa_2},\quad C_4=-\frac{\tau+t_2}{2}.
\]
Since $\Gamma(-3/2)=4\sqrt\pi/3$,
\begin{equation}\label{eq:normalization}
 \frac{C_3}{\Gamma(-3/2)}=A,\qquad
 \frac{[X^3]e^{C}}{\Gamma(-3/2)}=e^\nu A.
\end{equation}
This makes the factorial and Gamma normalization transparent.

For noninteger $\alpha$ the exact coefficient identity is
\[
 [z^n](1-z/\rho)^\alpha
 =\rho^{-n}\frac{\Gamma(n-\alpha)}{\Gamma(-\alpha)\Gamma(n+1)}.
\]
The fixed-order formal Gamma-ratio expansion follows from
\[
 \log\frac{\Gamma(n+a)}{\Gamma(n+b)}
 =(a-b)\log n+
 \sum_{k\geq1}\frac{(-1)^{k+1}\{B_{k+1}(a)-B_{k+1}(b)\}}
                         {k(k+1)n^k}.
\]
Here the infinite sum is used only after truncation to a fixed order.
For an independently coded alternative, expand $b(\tau+y)$ in ordinary
Taylor powers, invert $X^2=1-(1+y/\tau)e^{-(b(\tau+y)-b(\tau))}$,
and use the central-binomial expansion to handle half-integer Gamma
ratios. The companion implements this second route without Euler
derivative or Gaussian-moment inputs.

For example, $d=15/8-(5/2)C_5/C_3$. After division by $ve^{v\nu}$,
the $X^5$ coefficient of $e^{vC}$ is $C_5+vC_2C_3$, and its $X^7$
coefficient is
\[
 C_7+v(C_2C_5+C_3C_4)+\frac{v^2}2C_2^2C_3.
\]
The Gamma ratios $-5/2$ and $35/4$ then recover
\eqref{eq:d1v}--\eqref{eq:d2v}. To promote these formal local calculations
to a separate analytic coefficient proof one must supply suitable
$\Delta$-domain continuation and transfer remainder hypotheses, as in
the established subcritical theory \cite{subcritical}. That additional
route is not silently assumed here.

\section{Component law moments and connectivity}

Put
\begin{equation}\label{eq:aq}
 a=\frac{5\tau}{2},\qquad
 q=\tau d-\frac{35}8t_2+\frac52\tau^2.
\end{equation}
Dividing the marker expansion by its value at $1$ gives, uniformly on
fixed compact sets,
\begin{equation}\label{eq:pgffirst}
 P_n(v):=\EE v^{K_n}
 =ve^{\nu(v-1)}\left(1+\frac{a(v-1)}n+O(n^{-2})\right).
\end{equation}
In particular, the classical limit is $1+\Pois(\nu)$ \cite{bgkn}.
The next statement strengthens the approximation at finite $n$.

\begin{theorem}[Total variation refinement]\label{thm:TV}
For $\lambda_n=\nu+5\tau/(2n)$,
\[
 \TV\bigl(\mathcal L(K_n),\mathcal L(1+\Pois(\lambda_n))\bigr)
 =O(n^{-2}).
\]
\end{theorem}
\begin{proof}
The parameter is positive. Fix a real $r>1$. The comparison PGF is
$Q_n(v)=v\exp(\lambda_n(v-1))$. Equation~\eqref{eq:pgffirst}, proved
complex-uniformly on the circle $|v|=r$, gives
\[
 \sup_{|v|=r}|P_n(v)-Q_n(v)|\leq C_r n^{-2}
\]
for all sufficiently large $n$. Write $p_{n,k}=[v^k]P_n$ and
$q_{n,k}=[v^k]Q_n$. Cauchy's coefficient estimate gives
\[
 |p_{n,k}-q_{n,k}|\leq C_r n^{-2}r^{-k}\qquad(k\geq0).
\]
Both constant coefficients are zero for $n\geq1$. Therefore the entire
coefficient sum, including the Poisson tail above $n$, satisfies
\[
 \frac12\sum_{k\geq1}|p_{n,k}-q_{n,k}|
 \leq\frac{C_r}{2n^2}\sum_{k\geq1}r^{-k}
 =\frac{C_r}{2(r-1)n^2}.
\]
This is exactly the total variation distance.
\end{proof}

To obtain second-order moments, write
$d_2(v)=e+hv+kv^2$, where
\[
 e=d_2(0),\qquad h=\frac72\tau d-\frac{35}8t_2,
 \qquad k=\frac{35}8\tau^2.
\]
On a fixed sufficiently small disk about $v=1$, the normalized expansion
is nonzero for all sufficiently large $n$. Taking its analytic logarithm
gives
\begin{align}
 \log P_n(v)
 &=\log v+\nu(v-1)+\frac{a(v-1)}n \nonumber\\
 &\quad+\frac{q(v-1)+(5\tau^2/4)(v-1)^2}{n^2}+O(n^{-3}).
 \label{eq:logPGF}
\end{align}
Indeed the second-order polynomial before simplification is
$(h-da)(v-1)+(k-a^2/2)(v^2-1)$ and
$k-a^2/2=5\tau^2/4$.
The error is analytic and uniformly $O(n^{-3})$ on a slightly larger
disk, so Cauchy estimates justify differentiating it.

\begin{corollary}[Moments and connectivity]\label{cor:moments}
The component mean and variance satisfy
\begin{align}
 \EE K_n&=1+\nu+\frac an+\frac q{n^2}+O(n^{-3}),\label{eq:mean}\\
 \Var K_n&=\nu+\frac an+
           \frac{q+5\tau^2/2}{n^2}+O(n^{-3}).\label{eq:variance}
\end{align}
The probability of connectivity satisfies
\begin{equation}\label{eq:connectivity}
 \frac{c_n}{g_n}=e^{-\nu}\left[
 1-\frac{5\tau}{2n}
 +\frac{-\tau d+(35/8)t_2+(15/8)\tau^2}{n^2}
 +O(n^{-3})\right].
\end{equation}
\end{corollary}
\begin{proof}
The first derivative of $\log P_n$ at $1$ is the mean; its second
derivative plus its first derivative is the variance. Apply these
identities to \eqref{eq:logPGF}. Dividing \eqref{eq:conna} by
\eqref{eq:alla} gives the second coefficient
$-h+da+a^2-k=-\tau d+(35/8)t_2+(15/8)\tau^2$ in
\eqref{eq:connectivity}.
\end{proof}

Numerically,
\[
\begin{aligned}
 a&=0.42691246703157255852034948\ldots,\\
 q&=1.28213096120148270341548384\ldots,\\
 q+5\tau^2/2&=1.35503266300427611409382198\ldots,\\
 -\tau d+(35/8)t_2+(15/8)\tau^2
   &=-1.15455298304659423472839209\ldots,\\
 e^{-\nu}&=0.861666699427543\ldots.
\end{aligned}
\]
In particular, connectivity approaches its limiting probability from
below to first order. All higher fixed-order component and connectivity
corrections can be obtained by algebraic operations on the same marked
generator.

\section{Logarithmic expansion and a safe threshold inverse}

Take either $a_n=c_n$ or $a_n=g_n$. Write $A_*=A$ or $e^\nu A$,
respectively, and $e_j=d_j(0)$ or $d_j(1)$. Combining
Theorem~\ref{thm:main} with Stirling's expansion gives, for every fixed
integer $m\geq0$,
\begin{equation}\label{eq:logcount}
 \log a_n=F_m(n)+O(n^{-m-1}),
\end{equation}
where
\begin{equation}\label{eq:F}
 F_m(x)=x\left(\log\frac{x}{\rho}-1\right)-2\log x
       +\log(A_*\sqrt{2\pi})+\sum_{j=1}^m\frac{\ell_j}{x^j}.
\end{equation}
The $\ell_j$ are obtained by adding the Stirling series to the formal
logarithm of $1+\sum e_jx^{-j}$. In particular,
\begin{align*}
 \ell_1&=e_1+\frac1{12},\\
 \ell_2&=e_2-\frac12e_1^2,\\
 \ell_3&=e_3-e_1e_2+\frac13e_1^3-\frac1{360},\\
 \ell_4&=e_4-e_1e_3-\frac12e_2^2+e_1^2e_2-\frac14e_1^4.
\end{align*}
Thus every $\ell_j\in\QQ(\tau)$, for either sequence.

Both sequences are eventually strictly increasing. Adding the new label
as an isolated vertex injects the all-graph class on $[n]$ into that on
$[n+1]$, with other graphs outside its image. Attaching the new label as
a leaf to vertex $1$ similarly injects the connected class; for $n\geq2$
a tree with the new vertex of degree two is outside the image. Define
\[
 N(X)=\min\{n\geq1:a_n\geq X\}
\]
for sufficiently large real $X$.

Let $L=\log X$. A convenient explicit starting point for a continuous
inverse is
\begin{equation}\label{eq:Lambert}
 x_0=\frac{L}{W(L/(e\rho))},
\end{equation}
where $W$ is the positive real Lambert function. It solves
$x_0(\log(x_0/\rho)-1)=L$. For all large $L$, $F_m$ is strictly
increasing in the relevant range and has a unique large root $r_m$ of
$F_m(r_m)=L$. Since the lower-order residual at $x_0$ is $O(\log x_0)$
and the derivative is asymptotic to $\log x_0$, $r_m=x_0+O(1)$.
One bounded correction is also explicit. Set
$k_*=\log(A_*\sqrt{2\pi})$ and $\Lambda_0=\log(x_0/\rho)$. For every
fixed $m\geq0$,
\begin{equation}\label{eq:inverseshift}
 r_m=x_0+\frac{2\log x_0-k_*}{\Lambda_0}
           +O\!\left(\frac1{x_0\Lambda_0}\right).
\end{equation}
To verify this, the displayed displacement $h$ is bounded and tends to
$2$. Expanding $F_m(x_0+h)-L$ gives
$\Lambda_0h-2\log x_0+k_*+O(x_0^{-1})=O(x_0^{-1})$.
The curvature of the leading term, the change in $-2\log x$, and any
retained inverse-power terms are all $O(x_0^{-1})$. Division by
$F_m'(x)\sim\Lambda_0$ gives \eqref{eq:inverseshift}. Thus $r_m-x_0\to2$;
this is a statement about the continuous comparison root.

Newton iteration uses the explicit derivative
\begin{equation}\label{eq:Newton}
 F_m'(x)=\log(x/\rho)-\frac2x
             -\sum_{j=1}^m\frac{j\ell_j}{x^{j+1}}.
\end{equation}

\begin{theorem}[Integer-safe inverse]\label{thm:inverse}
For each fixed $m\geq0$ there is $K_m>0$ such that, for all sufficiently
large $X$, with $r=r_m(X)$ and
\[
 \delta=\frac{K_m r^{-m-1}}{\log r},
\]
one has
\begin{equation}\label{eq:inverse}
 \ceil{r-\delta}\ \leq\ N(X)\ \leq\ \ceil{r+\delta}.
\end{equation}
\end{theorem}
\begin{proof}
Write the error in \eqref{eq:logcount} as $E_m(n)$, with
$|E_m(n)|\leq C_mn^{-m-1}$ beyond a fixed integer threshold. Increase
that threshold so that for $|x-r|\leq2$,
\[
 F_m'(x)\geq\tfrac12\log r,
 \qquad |E_m(n)|\leq2C_m r^{-m-1}\quad(n\in\NN,\ |n-r|\leq2).
\]
Choose $K_m=8C_m$, or any larger positive constant; enlarge $C_m$ to be
positive if necessary. Eventually $\delta<1$. The explicit integer test
points are
\[
 n_-:=\ceil{r-\delta}-1,\qquad n_+:=\ceil{r+\delta}.
\]
Both lie within distance two of $r$, with $n_-<r-\delta$ and
$n_+\geq r+\delta$. The mean value theorem gives
\[
 F_m(n_-)<L-4C_mr^{-m-1},\qquad
 F_m(n_+)\geq L+4C_mr^{-m-1}.
\]
Allowing for $E_m$ shows $a_{n_-}<X$ and $a_{n_+}\geq X$.
Eventual monotonicity, together with $X$ exceeding the finitely many
earlier values, proves $n_-+1\leq N(X)\leq n_+$, which is
\eqref{eq:inverse}.
\end{proof}

The two ceilings agree except in a shrinking neighborhood of an integer.
One cannot replace the enclosure by an unconditional formula
$N(X)=\ceil{r_m(X)}$. Exact-threshold inputs $X=a_n$ are particularly
sensitive: even a tiny positive error $r_m(X)-n$ makes the latter ceiling
wrong. The companion can exhibit this numerically, but such diagnostics
are not proofs of a remainder constant or of a finite-input enclosure.
Theorem~\ref{thm:inverse} is an asymptotic existence result. A certified
finite-input algorithm would additionally require evaluated error
constants, a starting threshold, and interval evaluation of the root.

\section{Reproducibility and the meaning of the checks}

The companion is self-contained after its documented Python dependencies
are installed. No network access is needed for its tests. It contains
frozen OEIS b-files with source URLs, exact coefficient fixtures, portable
programs, bounded commands, and deterministic release construction.
The report and code use the same definitions, while several independent
routes check their implementation.

The checks are deliberately separated by what they establish:
\begin{enumerate}
\item Exact block coefficients are computed from the square-root identity
and independently from $S=1+uS+uS^2$. Integer Bell recurrences implement
\eqref{eq:connexact}, and the exponential formula computes $g_n$.
All 200 positive-index terms of each included OEIS b-file agree.
The connected index-zero convention is checked separately.
\item Literal enumeration visits each simple edge set through $n=5$.
It accepts an edge set if some cyclic ordering of all vertices has no
alternating-endpoint pair of edges, and tests connectivity directly.
This is the circular straight-line characterization of outerplanarity.
An edge set is counted only once, not once per successful ordering.
It reproduces $c_1,\ldots,c_5=1,1,4,37,602$ and
$g_1,\ldots,g_5=1,2,8,63,893$.
\item Exact arithmetic in $\QQ(\tau)$ implements the finite Gaussian
rule through four corrections and checks the marked identities
\eqref{eq:d1v}--\eqref{eq:d2v}, as well as the moment and connectivity
algebra. Stored rational-basis certificates are compared exactly.
\item A separately coded ordinary-Taylor implicit inversion and
central-binomial transfer checks the coefficient values at high
precision. This is independent formal coefficient verification; it
is not a second analytic $\Delta$-domain proof.
\item Frozen decimal fixtures check numerical regression. Explicit
failure predicates remain active under \texttt{python -O}. Generated
results are compared between normal and optimized runs, and manifests
check file integrity.
\end{enumerate}

These finite checks do not prove asymptotic convergence, bound the
constants hidden in $O(\cdot)$, or certify a finite-input inverse.
Count-ratio, moment, and inverse residual tables, when produced, are
diagnostics only. The analytic theorem rests on the exact identities
and the uniform contour proof. High working precision and agreement
between implementations do not turn approximate decimals into rigorous
intervals.

The precise commands, dependency versions, source URLs, fixture hashes,
and generated files are described in the companion's README and source
provenance file. Output directories must be new; an existing output is
rejected rather than silently overwritten. The release manifest records
the bytes actually delivered.

\section{Source normalization and attribution}

The leading outerplanar theorem in published BGKN
\cite[Theorem 5.1, pp.~2102--2103]{bgkn} uses the same unrooted labeled
normalization as \eqref{eq:alla}. It prints the all-graph amplitude as
$0.018216$. Equations \eqref{eq:constants} and \eqref{eq:normalization}
give instead
\[
 e^\nu A=0.008096047476080704996204256500405217\ldots,
\]
while
\[
 \frac94e^\nu A=0.018216106821181586241459577125911738\ldots.
\]
Thus the printed amplitude agrees, to its printed precision, with
$9/4$ times the value obtained from the defining generating functions.
No cause of that discrepancy is established here.

Version distinctions matter. The earlier full preprint
\cite{bgknpreprint} prints $0.017657$, a different decimal.
The published component theorem already gives the compatible approximate
parameter $0.14889$ and connectivity probability $0.86166$; its
introductory decimals $0.14840$ and $0.86208$ are stale. The published
introduction also retains an inconsistent $n^{-3/2}$ outerplanar
exponent, while its main theorem has $n^{-5/2}$. These introductory
issues do not alter the correct exponent in the theorem.

Kang's thesis \cite[p.~55]{kang} already gives the correct analytic
connected-amplitude expression and the relation
$\alpha_0=\alpha_1\exp(C(R))$. Its printed all-graph decimal
$0.008095$ is inaccurate, not just a lower-precision rounding of
$0.008096047\ldots$. The correct analytic normalization was therefore
already present in the source set. This report resolves numerical
normalization reproducibly; it does not announce a new leading constant.

The bounded source pass included the published and preprint BGKN
versions, Kang's thesis, the labeled section of the Bodirsky--Kang
sampling paper, the subcritical graph-class framework, and Finch's
outerplanar growth-constant note \cite{finch}. Finch's block convention
excludes $K_2$ from his block function and compensates with an extra
linear term in the rooted equation; combining them yields
\eqref{eq:block}. Unlabelled outerplanar graph work \cite{unlabelled}
concerns different sequences and constants.

No explicit first-four correction table, two-order component-moment
formula, or threshold inverse for these two labeled sequences was found
in the inspected source portions and bounded searches. This is a source
coverage statement, not proof of worldwide novelty. The leading theorem,
the shifted-Poisson limit, and the general machinery remain prior work.

\section{Further questions}

Several next steps would add information beyond merely printing more
terms from the same generator.
\begin{enumerate}
\item Evaluate explicit remainder constants and starting thresholds,
then combine them with interval arithmetic to turn
Theorem~\ref{thm:inverse} into a finite-input certified algorithm.
\item Determine the growth of $d_j(v)$ as $j$ increases. The fixed-order
proof does not address optimal truncation, divergence, or exponentially
small corrections.
\item Extract higher-order signed corrections to the shifted-Poisson
law and quantify their total-variation effect. Positivity of a proposed
finite-order distributional approximation requires care.
\item Extend the marked exact cancellation to edge statistics jointly
with components, retaining complex uniformity and the abstract-graph
normalization. The leading edge-count limit theory is already classical.
\end{enumerate}

\clearpage
\begin{thebibliography}{10}
\bibitem{oeisconn} OEIS Foundation Inc.,
\emph{A097998: Number of connected outerplanar graphs on $n$ labeled nodes},
\url{https://oeis.org/A097998}. Positive-index b-file checked 3 October 2026.
\bibitem{oeisall} OEIS Foundation Inc.,
\emph{A098000: Number of outerplanar graphs on $n$ labeled nodes},
\url{https://oeis.org/A098000}. Positive-index b-file checked 3 October 2026.
\bibitem{bk} M. Bodirsky and M. Kang,
\emph{Generating Outerplanar Graphs Uniformly at Random},
Combinatorics, Probability and Computing 15 (2006), 333--343.
\url{https://doi.org/10.1017/S0963548305007303}.
\bibitem{bgkn} M. Bodirsky, O. Gim\'enez, M. Kang, and M. Noy,
\emph{Enumeration and limit laws of series-parallel graphs},
European Journal of Combinatorics 28 (2007), 2091--2105.
Published author-hosted version:
\url{https://web.mat.upc.edu/marc.noy/uploads/2013/05/Graphs-SP.pdf}.
\bibitem{bgknpreprint} M. Bodirsky, O. Gim\'enez, M. Kang, and M. Noy,
\emph{Enumeration and limit laws of series-parallel graphs},
earliest full preprint version, arXiv:math/0512435.
\url{https://arxiv.org/abs/math/0512435}.
\bibitem{kang} M. Kang,
\emph{Random Planar Structures and Random Graph Processes},
habilitation thesis, Humboldt-Universit\"at zu Berlin, June 2007,
Chapter 5, especially p.~55.
\href{https://edoc.hu-berlin.de/server/api/core/bitstreams/23f83814-f327-4458-9787-474e750b8262/content}{University repository PDF}.
\bibitem{subcritical} M. Drmota, \'E. Fusy, M. Kang, V. Kraus, and J. Ru\'e,
\emph{Asymptotic study of subcritical graph classes},
SIAM Journal on Discrete Mathematics 25 (2011), 1615--1651.
\url{https://arxiv.org/abs/1003.4699}; author-hosted version:
\url{https://web.mat.upc.edu/juan.jose.rue/Research/Subcritical.pdf}.
\bibitem{finch} S. R. Finch, \emph{Planar graph growth constants},
OEIS-hosted author note,
\url{https://oeis.org/A097999/a097999.pdf}.
\bibitem{unlabelled} M. Bodirsky, \'E. Fusy, M. Kang, and S. Vigerske,
\emph{Enumeration and asymptotic properties of unlabeled outerplanar graphs},
\url{https://arxiv.org/abs/math/0511422}.
\end{thebibliography}
\end{document}
